\VerdierPhysicalPageAnchor{0081}%
\VerdierSemanticAnchor{ch1-2-5-11}
\noindent\textbf{2.5.11.}\quad
\VerdierSemanticAnchor{ch1-2-5-11-homotopism}
$X$ と $Y$ を $C$ の二つの複体とする。$\operatorname{comp}(C)$ の射
$f\colon X\longrightarrow Y$ は、その $K(C)$ における像 $\dot f$ が
同型であるとき、\emph{ホモトピー同値}と呼ばれる。
\VerdierSemanticAnchor{ch1-2-5-11-homotopic-complexes}
複体 $X$ と $Y$ は、$K(C)$ において同型であるとき、\emph{ホモトピー
同値}であるという。

\VerdierSemanticAnchor{ch1-2-5-12}
\noindent\textbf{2.5.12.}\quad
\VerdierSemanticAnchor{ch1-2-5-12-setup}
$X,Y$ を $C$ の二つの複体、
$f,g\colon X\rightrightarrows Y$ を $\operatorname{comp}(C)$ の二つの射、
$a,b\colon f\rightsquigarrow g$ を $f$ と $g$ を結ぶ二つのホモトピーとする。
\VerdierSemanticAnchor{ch1-2-5-12-second-order-homotopy}
$a$ と $b$ を結ぶ\emph{二階ホモトピー}とは、
$\alpha\colon a\rightsquigarrow b$ と表される射の族
\[
  \alpha=(\alpha^i\colon X^{i+2}\longrightarrow Y^i,\ i\in\mathbb Z)
\]
であって、次を満たすものをいう：

\VerdierSemanticAnchor{ch1-2-5-12-eq1}
\begin{equation}
  b^i-a^i
  =\alpha^i d_X^{i+1}-d_Y^{i-1}\alpha^{i-1}.
  \tag{2.5.12.1}
\end{equation}

\VerdierSemanticAnchor{ch1-2-5-13}
\noindent\textbf{2.5.13.}\quad
最後に、同型 (2.4.1)

\VerdierSemanticAnchor{ch1-2-5-13-complex-opposite-isomorphism}
\begin{equation*}
  \iota_1\colon
  \operatorname{comp}(C)^\circ
  \longrightarrow
  \operatorname{comp}(C^\circ)
\end{equation*}
は、(2.5.7) により一意な同型

\VerdierSemanticAnchor{ch1-2-5-13-homotopy-opposite-isomorphism}
\begin{equation*}
  \iota_1\colon
  K(C)^\circ
  \longrightarrow
  K(C^\circ)
\end{equation*}
\VerdierSemanticAnchor{ch1-2-5-13-z-category-isomorphism}
を定め、これは $\mathbb Z$-圏の同型であることを注意しておく。
\VerdierSemanticAnchor{ch1-2-5-13-identification}
以下では、この同型によって圏 $K(C)^\circ$ と $K(C^\circ)$ を同一視することが多い。
